\chapter{Phonology}

\section{Vowels}

\subsection{Monophthongs}

Niklorsk has nine phonemic vowels, as shown in Table \ref{tab:vowel-inventory} below. In addition /ə/ is a marginal phoneme, occurring only in unstressed coda syllables as an allophone of /e/.

\begin{table}[!h]
    \begin{tabularx}{\textwidth}{X X X X}
        & Front & Central & Back \\
        Close & i, y & & ʉ, u \\
        Mid & e, œ & (ə) & o \\
        Open & æ & & ɑ \\
    \end{tabularx}
    \label{tab:vowel-inventory}
    \caption{Vowel inventory of Niklorsk}
\end{table}

Each vowel occurs in both long and short forms, as shown in Table \ref{tab:vowel-length}. In general, long vowels are tenser than their short counterparts.

\begin{table}[!h]
    \begin{tabularx}{\textwidth}{X X X X}
        & Front & Central & Back \\
        Close & ɪ iː, ʏ yː & & ɵ ʉː, ʊ uː \\
        Mid & ɛ eː, œ øː & (ə) & ɔ oː \\
        Open & æ æː & & ɑ ɑː \\
    \end{tabularx}
    \label{tab:vowel-length}
    \caption{Vowel length in Niklorsk}
\end{table}

Vowel length is marked in the orthography by doubling the consonant following the vowel. For example, \nikpp{mod}{muːð}{drink} and \nikpp{modd}{mʊðː}{meadow}. The following table shows the orthographic representation of the vowels. The low vowels /ɑː/ and /æː/ are raised to [ɔː] and [ɛː], respectively. (In some dialects, this vowel raising also affects short /æ/, so that /æ/ and /ɛ/ are merged.)

\begin{table}[!h]
    \begin{tabularx}{\textwidth}{X X}
        Orthography & Phoneme \\
        a & /ɑ/ \\
        e & /e/ \\
        i & /i/ \\
        o & /u/ \\
        u & /ʉ/ \\
        y & /y/ \\
        ä & /æ/ \\
        ö & /œ/ \\
        å & /o/ \\
    \end{tabularx}
    \label{tab:vowel-orthography}
    \caption{Orthographic representation of vowels in Niklorsk}
\end{table}

\subsection{Diphthongs}

\begin{table}
    \small
    \begin{tblr}{
        colspec={XX}
    }
        æ /aɛ̯/ & \nikpp{bærg}{baɛ̯rj}{heaven,} \\
        ee /eɪ̯~eːɪ̯/ &  \nikpp{seete}{seɪ̯tə}{lake,} \\
        au /œʏ̯/ & \\
    \end{tblr}
    \caption{Diphthongs in Niklorsk}
    \label{tab:diphthongs}
\end{table}

Niklorsk has three phonemic diphthongs, as shown in Table \ref{tab:diphthongs}. For the sake of prosody, diphthongs are treated as long vowels.


\section{Consonants}

The consonant inventory of Niklorsk is shown in Table \ref{tab:consonant-inventory} below. Vowels appear as either short or long, with the latter marked by a doubled consonant following the vowel.

\begin{table}
    \small
    \begin{tblr}{
        colspec={XXXXXX}
    }
        & Labial & Dental/ Alveolar & Palatal & Velar & Glottal \\
        Nasal       & m     &  n    &       &       & ŋ  \\
        Plosive     & p b   &  t d  &       & k g   & \\
        Fricative   & f ʋ   &  s    & ɕ     & ç     & h \\
        Liquid      &       &  l r   & j      &       & \\
       
    \end{tblr}
    \caption{Consonant inventory of Niklorsk}
    \label{tab:consonant-inventory}
\end{table}

\subsection{Plosives}
Voiceless plosives /p t k/ are aspirated when they occur at the onset of a stressed syllable, with /k/ often realized as [cʰ] instead of [kʰ]. For example, \nikpp{tor}{tʰuːɐ̯}{to believe,} \nikpp{kapp}{cʰɑpː}{coat,} \nikpp{pær}{pʰaɛ̯}{pen,} \nikpp{patat}{pɑːtʰɔːt}{potato.} This aspiration does not occur in unstressed syllables, at syllable codas, or after /s/ (with /sk/ becoming [ɕ]), e.g., \nikpp{strom}{struːm}{already.} Before front vowels in a stressed syllable, /t/ and /d/ are affricated to /ts/ and /dz/, respectively, e.g., \nikpp{titt}{t͡sɪtː}{to hide,} \nikpp{dyr}{d͡zyːɐ̯}{bear.} Long unvoiced plosives in intervocalic positions are voiced, i.e., /tː/ becomes [ðː], /kː/ becomes [gː], and /pː/ becomes [bː], e.g., \nikpp{atta}{ɑðːɔː}{into the.}

Most stops undergo lenition in the coda. At the end of a syllable /k/ becomes [c], /d/ becomes [ð], and /g/ becomes silent, e.g., \nikpp{tadd}{tʰɑðː}{father,} \nikpp{tak}{tʰɔːc}{brow,} \nikpp{tog}{tʰoː}{desire.}

Before front vowels, /g/ becomes [j], e.g., \nikpp{merge}{mɛɐ̯ɪ}{shirt,} \nikpp{Germanska}{jɛɐ̯mɑnʂkɔː}{German}. Similarly, Consonant clusters ending in /g/ and found in the syllable coda (excepting declension and other morphological endings) see the /g/ become [j], e.g., \nikpp{bærg}{baɛ̯rj}{heaven.} The unvoiced /k/ becomes [ɕ] before front vowels, e.g., \nikpp{käll}{ɕælː}{who.}

\subsection{Fricatives}

The fricative /s/ becomes [ʂ] (or sometimes [ɕ]) after /l/ or /r/, e.g., \nikpp{skers}{ɕɛɐ̯ʂ}{to sing (inf.)}, or before /t/ or /d/, e.g., \nikpp{forst}{fuːɐ̯ʂt}{time.} Final /ʋ/ is often simplified to [w], e.g., \nikpp{bliv}{bliːw}{man.} The sequence /st/ is written as \nk{sst} and pronounced [sːt], e.g., \nikpp{lasst}{lɑsːt}{salt.} This sequence cannot occur word-initially.

\subsection{Liquids}

Onset /r/ by itself or in a consonant cluster is realized as the uvular fricative [ʁ], e.g., \nikpp{brörd}{bʁøːɐ̯}{room,} \nikpp{ræn}{ʁaɛ̯n}{sweat.}

Syllable-final /r/ is derhotacized to form a /ɐ̯/ diphthong with the preceding vowel, maintaining the vowel length of the latter. Thus we have \nikpp{hör}{høːɐ̯}{pen,} \nikpp{börrs}{bœɐ̯ʂ}{forest.} If the preceding vowel is a diphthong, the /r/ is altogether dropped if /r/ is not part of a consonant cluster: \nikpp{bærg}{baɛ̯rj}{heaven} but \nikpp{seer}{seːɪ̯}{water.}

\section{Orthography}

\subsection{General rules}

Spelling and pronunciation are generally regular in Niklorsk, following the phonetic rules outlined above. The following provides to the orthographic representation of the consonants and vowels, as well as exceptions to the general rules.


\begin{enumerate}[wide, nosep]
    \item Final /d/ is silent after the liquids /l/ and /r/, the nasals /m/, /n/, and /ŋ/, and the fricatives /s/ and /ʋ/. This lenition results in the lengthening of the preceding consonant except for /s/ and /ʋ/ and the liquid /r/ which is derhotacized. The sequence \nk{gnd} is simplified to [ŋː] or just [ŋ] (cf. \ref{item:gn} below). This lenition is maintained when adding the weak suffixes, i.e., those starting with \nk{a}, \nk{e} or \nk{s}.
    \begin{itemize}
        \item \nikpp{börd}{bøːɐ̯}{eat (pres.)} but \nikpp{bördutt}{bøːɐ̯ðɵtː}{eat (prog.)}
        \item \nikpp{ngöld}{ŋœlː}{cow} and \nikpp{ngölddar}{ŋœlːaɐ̯}{cows}
        \item \nikpp{bängds}{bæŋːs}{to cover (inf.)} but \nikpp{bängdutt}{bænðɵtː}{cover (prog.)}
    \end{itemize}
    \item The digraph \nk{hj} represents /ç/
    \item The fricative /ɕ/ is written as \nk{sk} before front vowels or word-finally and as \nk{skj} elsewhere. If following /r/, however, /ɕ/ becomes [ʂ].
    \begin{itemize}
        \item \nikpp{niklorsk}{nɪˈkluːɐ̯ʂ}{Niklorsk} but \nikpp{mask}{mɔːɕ}{vermin}
    \end{itemize}
    \item \nk{-ans} is usually [ɔns] except notably in French loanwords where the original nasal sound is approximated as [ɔŋs].
    \begin{itemize}
        \item \nikpp{ambulans}{ɑmbʰuːlɔŋs}{ambulance}
        \item \nikpp{dans}{dɔŋs}{dance}
    \end{itemize}
    \item The final /n/ in the plural suffix \nk{-(e)n} as well as the passive suffix \nk{-(e)n} is silent.
    \begin{itemize}
        \item \nikpp{aften}{ɑftə}{books} same as \nikpp{afte}{ɑftə}{book}
    \end{itemize}
    \item The sequence \nk{gn} represent /ŋn/ (often realized as [ŋːn]) and can only appear in the syllable coda or followed by a vowel or either of /s/ or /t/. In the latter case, as well as in the syllable coda, the /n/ is silent.\label{item:gn}
    \begin{itemize}
        \item \nikpp{begns}{bɛŋːs}{read (inf.)} but \nikpp{begnar}{bɛŋːnɔːɐ̯}{read (past)}
    \end{itemize}
    \item The sequence \nk{-tion}, \nk{-sion} and \nk{-son} in Latinate loanwords are both pronounced [ʂuːn]. Some speakers, mainly due to Swedish influence, pronounce it as [huːn].
    \begin{itemize}
        \item \nikpp{information}{ɪnfʊɐ̯mɔːʂu:n ɪnfʊɐ̯mɔːhu:n}{information}
        \item \nikpp{situasion}{sɪtʉaʂu:n sɪtʉahu:n}{situation}
        \item \nikpp{version}{vɛɐ̯ʂu:n vɛɐ̯hu:n}{version}
        \item \nikpp{person}{pɛɐ̯ʂu:n pɛɐ̯hu:n}{person}
    \end{itemize}
    \item In colloquial speech, the past participle ending \nk{-agne} is pronounced [ɑŋ] (or by some speakers [ɑŋː]) instead of [ɑŋːnə]. Similarly, the present participle ending \nk{-ade} is pronounced [ɔː] instead of [ɔːðə].
    \item The sequence \nk{kv}, \nk{gv} and \nk{hv} are all realized as [ɕʋ] or [ɕv].
    \begin{itemize}
        \item \nikpp{kvæld}{ɕʋɛlː}{evening}
        \item \nikpp{i kvett}{ɪ ɕʋɛtː}{yesterday}
        \item \nikpp{gvers}{ɕʋɛɐ̯ʂ}{sky}
    \end{itemize}
    \item Nasal /m/, /n/ and /ŋ/ assimilate to the place of articulation of the following consonant. 
    \item In unstressed syllables, the sequence \nk{ge} if not followed by a consonant in the same syllable is pronounced [ɪ] by most speakers instead of the standard [jə].
    \begin{itemize}
        \item \nikpp{merge}{mɛɐ̯ɪ}{shirt}
        \item \nikpp{Norge}{nuɐ̯ɪ}{Norway}
        \item \nikpp{tage}{tɔːɪ}{to bring} but \nikpp{tagett}{tɔːjətː}{to bring (prog.)}
    \end{itemize}
    \item 
\end{enumerate}

\subsection{Pronunciation of common words}

Some common words have irregular pronunciations. This section provides a survey of such words.